\section{Results}\label{sec:results}

\subsection{All Public CPU Comparators vs.\ \ultra{}}\label{sec:compare}

\Cref{fig:money} draws live +LM at 160\,ms and the utterance-end commit pass against every timed
public comparator. \Cref{tab:measured} keeps those two rows plus the 320\,ms and greedy modes of the
same 60\,MB INT8 artifact, beside all \ColdCompleteCompCount{} public open-weights CPU comparators
with complete 730-utterance LibriSpeech DEV panel WER. \ColdHostSentence{} \ColdSubsetSentence{} \ColdOrderSentence{}

\para{Speed and lowest-power operating points.} \ColdHeadSentence{} On the exact same weights, matching
Zipformer INT8's 320\,ms chunk length (\ultra{} live +LM, 320\,ms) yields our lowest active-core energy
and best live WER at \VsWideSpeed$\times$ real time, \VsWideWatts\,J per audio-second, and \VsWideWer\%
panel WER (\CandLiveLmEightFullWer\% full DEV)---\VsWideSpeedX$\times$ Zipformer's paired throughput at
\VsWideEnergyPct\% of its active core energy. Greedy 160\,ms live decoding reaches \VsGreedySpeed$\times$
real time in \ColdRcTwoLiveGreedyRss\,MB RAM at \VsGreedyWatts\,J per audio-second (\VsGreedyWer\% panel
WER). At utterance endpoints, the commit pass runs at \ColdFinalCorePct\% of one quiet Zen~2 core
(\ColdFinalVsZipSpeed{}).

\para{Accuracy.} \ColdAccuracySentence{} Notably, while \ultra{} uses Libri-Light pseudo-labels
generated by the Zipformer comparator (R0), its utterance-segmented commit pass matches streaming
Zipformer INT8's pooled panel WER (\ColdHeadWer\% vs.\ \ColdZipWer\%, \VsVerdictTwo{} it under the
paired test of \cref{sec:versus}), with lower dev-other WER (\ColdHeadWerOther\% vs.\ \ColdZipWerOther\%)
and equal dev-clean WER (\ColdHeadWerClean\% vs.\ \ColdZipWerClean\%). Larger 600M-parameter models
achieve lower WER at much higher CPU and energy cost (offline Parakeet TDT v3 on quiet H1; Nemotron
560\,ms on the quiet 73-clip timing subset; Parakeet unified stream 160\,ms on Xeon~8358).
\RelSubsetNote{}

\para{Memory and active-core energy.} \ColdMemorySentence{} \ColdWattsSentence{} The bottom panel of
\cref{fig:money} draws only the matched R6 rows: live +LM is \ColdLiveWatts\,J per audio-second,
\VsEnergyPctOne\% of Zipformer INT8 (\ColdZipRsixWatts\,J) on the same blocks, with the idle floor
subtracted. The commit pass is \ColdHeadWatts\,J per audio-second. Moonshine, Parakeet, and Nemotron
joules in \cref{tab:measured} include the core idle floor and are not a paired ratio. The 320\,ms and
greedy rows in \cref{tab:measured} are \VsWideWatts\,J and \VsGreedyWatts\,J.

\subsection{Head-to-Head with Zipformer INT8}\label{sec:versus}

Zipformer INT8 is the strongest public streaming comparator near our cost point and the teacher of our
pseudo-labels, so we compare against it directly on the same utterances and the same core
(\cref{tab:versus}), timed in the same three interleaved blocks on quiet H1; speed, energy and memory
ratios are medians of the within-block ratios. WER differences are paired with speaker-cluster
bootstrap. Because the 730-utterance panel overlaps the selection half of DEV, we always report the
unselected confirmation half (C, \CandCUtts{} utterances; \cref{sec:candidate}) beside it: on C, live
+LM trails Zipformer by \CandLiveLmEightCDelta\pp (320\,ms, bounds \CandLiveLmEightCLo{} to
\CandLiveLmEightCHi) and \CandLiveLmCDelta\pp (160\,ms, bounds \CandLiveLmCLo{} to \CandLiveLmCHi).
The 160\,ms live stream uses \VsCpuPctOne\% of Zipformer's CPU time (\VsSpeedXOne$\times$ throughput),
\VsEnergyPctOne\% of its active core energy, and \VsRssPctOne\% of its memory with half the chunk
(\VsChunkOne\,ms vs.\ \VsZipChunk\,ms); its panel WER is \VsVerdictOne{} Zipformer's (\VsDeltaOne\pp,
bounds \VsLoOne{} to \VsHiOne\pp). The final commit pass on the same weights is \VsVerdictTwo{}
Zipformer's (\VsDeltaTwo\pp, bounds \VsLoTwo{} to \VsHiTwo\pp; \CandFinalCDelta\pp on C) at
\VsSpeedXTwo$\times$ its single-core speed and \VsEnergyPctTwo\% of its active core energy.

\begin{table}[t]
\centering
\caption{\textbf{Head-to-head with Zipformer INT8}: same 730 DEV utterances (the selection panel; the unselected-half gaps are in \cref{sec:candidate}), same quiet Zen~2 core, same interleaved blocks, median of three. Ratios are medians of within-block ratios. $\Delta$ is paired (speaker-cluster bootstrap); negative means fewer errors. \textsuperscript{a}Live stream plus a commit pass at reference utterance boundaries.}
\label{tab:versus}
\footnotesize
\setlength{\tabcolsep}{1.5pt}
\begin{tabular}{@{}lrrr@{}}
\toprule
Metric & Zipformer & \ultra{} live & \ultra{} final \\
\midrule
Live chunk (ms) & 320 & 160 & 160\textsuperscript{a} \\
Pooled WER (\%) & 5.25 & 5.54 & 5.00 \\
$\Delta$ WER vs.\ Zip.\ (pp) & --- & +0.28 & $-$0.26 \\
\quad one-sided 95\% bounds & --- & [$-$0.02, +0.58] & [$-$0.57, +0.05] \\
1-core speed ($\times$ RT) & 14.0 & 32.7 & 16.8 \\
J per audio-s (core) & 0.573 & 0.328 & 0.661 \\
Peak RSS (MB) & 152 & 109 & 387 \\
\bottomrule
\end{tabular}
\end{table}


\subsection{Registered Look and Cross-ISA Reproducibility}\label{sec:candidate}

The \ReleaseVer{} acoustic model (RC2, identical weights in v3.2.0) was trained on \CandTrainKh\,h of
logged audio (adding Libri-Light large to the lineage) and frozen before one registered look on the
panel and the unselected C half (\CandCUtts{} utterances; \cref{tab:candidate}). On the panel, the
160\,ms live +LM stream trails Zipformer by \CandLiveLmPanelDelta\pp (upper \CandLiveLmPanelHi\pp), so
non-inferiority at the registered \CandNiMargin\pp margin is not shown (Holm-adjusted
$p=\CandHolmNi$); on C the gap is \CandLiveLmCDelta\pp (320\,ms: \CandLiveLmEightCDelta\pp). The final
pass reaches \CandFinalWer\% on the panel (\CandFinalCDelta\pp vs.\ Zipformer on C; Holm-adjusted
$p=\CandHolmFinal$). RC2 made \CandFinalFullEdits{} full-DEV final edits against a registered adoption
gate of \CandGateFinalFull{} and missed both panel edit targets (\CandGateLivePanel{} live,
\CandGateFinalPanel{} final). \ultra{} \ReleaseVer{}/v3.2.0 is released by owner decision with these
misses stated explicitly.

\ultra{} preserves exact cross-architecture reproducibility: \CandIsaRuns{} witness decodes on Intel
Xeon (AVX2 and AVX-512 VNNI) and Apple M3 (NEON) across live greedy, both live LM policies, and the
final pass reproduce the Zen~2 transcripts on all \ScalarUtts{} panel utterances. On main from
\ScalarMainCommit{} (including v3.2.0), a forced scalar reference backend also gives
\ScalarDiffs/\ScalarUtts{} differing utterances across AVX2, AVX-512 VNNI, NEON, and scalar.

\begin{table*}[t]
\centering
\caption{\textbf{RC2 (released as v3.1.0), single registered look.} It missed its registered gates and is released by owner decision, not as a gate pass. $\Delta$ is RC2 minus Zipformer INT8 (or minus shipped v3.0.0) on the same utterances, with one-sided 95\% bounds from the pinned scorer; negative means fewer errors. The 730-utterance panel overlaps the half of DEV used for selection; the C half (2,785 utterances) was never used for selection. Full-DEV WER is descriptive. Live rows compare with shipped v3.0.0 live greedy (160\,ms); the final rows with its final pass. Final rows run the 160\,ms live stream and commit with a 640\,ms chunk at reference utterance boundaries.}
\label{tab:candidate}
\footnotesize
\setlength{\tabcolsep}{4pt}
\begin{tabular}{@{}lrrrlll@{}}
\toprule
Mode & Chunk (ms) & Panel WER & Full-DEV WER & $\Delta$ vs.\ Zip.\ (panel) & $\Delta$ vs.\ Zip.\ (C half) & $\Delta$ vs.\ v3.0.0 (panel) \\
\midrule
Live, greedy & 160 & 6.23 & 6.83 & +0.98 [+0.66, +1.29] & +1.59 [+1.31, +1.90] & $-$0.62 [$-$0.86, $-$0.38] \\
Live + 34.5\,MiB LM & 160 & 5.54 & 6.10 & +0.28 [$-$0.02, +0.58] & +0.81 [+0.62, +1.02] & $-$1.31 [$-$1.63, $-$1.00] \\
Live + 34.5\,MiB LM & 320 & 5.50 & 5.87 & +0.24 [$-$0.04, +0.52] & +0.65 [+0.47, +0.84] & $-$1.35 [$-$1.68, $-$1.03] \\
Final, M5 LM & 640 & 5.00 & 5.22 & $-$0.26 [$-$0.57, +0.05] & $-$0.11 [$-$0.35, +0.11] & $-$0.06 [$-$0.28, +0.17] \\
Final, 146.8\,MiB LM & 640 & 5.09 & 5.24 & $-$0.17 [$-$0.49, +0.15] & $-$0.10 [$-$0.35, +0.12] & --- \\
\midrule
Zipformer INT8 & 320 & 5.25 & 5.37 & --- & WER 5.13 & --- \\
\bottomrule
\end{tabular}
\end{table*}


\subsection{Quiet-Speech Level Normalisation (v3.2.0 \texttt{-{}-rms-norm})}\label{sec:level}

The runtime-only v3.2.0 release (build commit \VThreeTwoBuildCommit{}, signed tag \VThreeTwoTagObject{},
keeping identical \ReleaseVer{} RC2 weights) adds an opt-in final-pass RMS level normaliser
(\texttt{-{}-rms-norm}, off by default). On the unselected DEV confirmation half (C) with the frozen M5
commit pass, \texttt{-{}-rms-norm} left native 0\,dB accuracy unchanged within noise
(\AudioZOneCleanDelta\pp, one-sided 95\% upper \AudioZOneCleanUpper\pp) while improving synthetic
$-$20\,dB and $-$30\,dB attenuated speech by \AudioZOneTwentyDelta\pp (upper \AudioZOneTwentyUpper\pp)
and \AudioZOneThirtyDelta\pp (upper \AudioZOneThirtyUpper\pp). On automatic-endpoint chapter streams
(DEV selection half S, 90 chapters, verifying hangover, EOF, deadline, and PCM-cap triggers), it changed
pooled WER by \AudioZOneAutoCleanDelta\pp at 0\,dB, \AudioZOneAutoTwentyDelta\pp at $-$20\,dB, and
\AudioZOneAutoThirtyDelta\pp at $-$30\,dB (uppers \AudioZOneAutoCleanUpper{},
\AudioZOneAutoTwentyUpper{}, and \AudioZOneAutoThirtyUpper\pp). Paired pinned-core H0 relative CPU
overhead stayed within \AudioZOneCpuEngineUpper\% for the commit binary and \AudioZOneCpuStreamUpper\%
for stream commit mode, with \AudioZOneCrossIsaDiffs/\AudioZOneCrossIsaClips{} committed-text
differences across AVX2, scalar, and NEON on the \AudioZOneCrossIsaCommit{} cross-ISA witness builds
(\AudioZOneOracleTupleDiffs/\AudioZOneCrossIsaClips{} full-tuple differences vs.\ the independent
oracle with flag on, \AudioZOneCrossIsaDiffs{} committed-text differences).

\subsection{Full-DEV Accuracy and INT8 Parity}\label{sec:fulldev}

\Cref{tab:fulldev} reports full LibriSpeech DEV (\RelFullUtts{} utterances: \CandFullCleanUtts{}
dev-clean and \CandFullOtherUtts{} dev-other) accuracy for all four modes of our single 60\,MB INT8
artifact under the same unified scorer. Full DEV contains the selection half, so these numbers are
descriptive. On the prior release artifact of identical architecture, prepared INT8 matched FP32 PyTorch
across all \RelFullUtts{} utterances within \RelIntEightDelta\pp pooled WER (one-sided 95\%
speaker-cluster bootstrap bounds \RelIntEightLo{} to \RelIntEightHi\pp), well inside the pre-registered
parity gate.

\begin{table}[t]
\centering
\caption{\textbf{Full LibriSpeech DEV (\RelFullUtts{} utterances)} WER (\%) across all four operating
points of our single 60\,MB INT8 artifact (unified \texttt{ultra-score}).}
\label{tab:fulldev}
\footnotesize
\setlength{\tabcolsep}{3pt}
\begin{tabular}{@{}lrrr@{}}
\toprule
Mode (same INT8 weights) & dev-clean & dev-other & Pooled \\
\midrule
Live greedy (160\,ms, $c=4$) & \CandLiveGreedyFullClean & \CandLiveGreedyFullOther & \CandLiveGreedyFullWer \\
Live +LM (160\,ms, $c=4$) & \CandLiveLmFullClean & \CandLiveLmFullOther & \CandLiveLmFullWer \\
Live +LM (320\,ms, $c=8$) & \CandLiveLmEightFullClean & \CandLiveLmEightFullOther & \CandLiveLmEightFullWer \\
\textbf{Final commit (640\,ms, M5)} & \textbf{\CandFinalFullClean} & \textbf{\CandFinalFullOther} & \textbf{\CandFinalFullWer} \\
\bottomrule
\end{tabular}
\end{table}

\subsection{Paced Always-On Live Capacity on One Quiet Zen~2 Core}\label{sec:capacity}

To measure real-time always-on capacity and stable-word lag under live packet pacing, we run
$N\in\{1,4,8,16\}$ concurrent live streams (\texttt{ultra stream}, real-time pacing, 20\,ms packets,
per-packet wakes) pinned to a single quiet Zen~2 core on the prior release artifact of identical
compute cost (\cref{tab:capacity}). Sixteen simultaneous live streams occupy \RelCapSixteenBusy\% of one
core at \RelCapSixteenNetW\,W net (\RelCapSixteenGrossW\,W gross) core power per stream and stable-word
lag \RelCapSixteenLagPfifty/\RelCapSixteenLagPninetyFive\,ms (p50/p95). A lone paced stream costs more
per stream (\RelCapOneNetW\,W net, \RelCapOneBusy\% busy) because it wakes an otherwise idle core at the
packet rate and amortises model load over a single 128\,s stream.

\begin{table}[t]
\centering
\caption{\textbf{Always-on live capacity on one quiet Zen~2 core}
(real-time pacing, 20\,ms packets, continuous 128\,s stream of 22 DEV clips; prior release artifact of
identical per-frame cost). Net W subtracts the idle core floor; gross W is total core power divided by $N$.}
\label{tab:capacity}
\footnotesize
\setlength{\tabcolsep}{3.5pt}
\begin{tabular}{@{}rrrr@{}}
\toprule
Streams ($N$) & Core busy (\%) & Net / gross W & Lag p50 / p95 (ms) \\
\midrule
1 & \RelCapOneBusy & \RelCapOneNetW{} / \RelCapOneGrossW & \RelCapOneLagPfifty{} / \RelCapOneLagPninetyFive \\
4 & \RelCapFourBusy & \RelCapFourNetW{} / \RelCapFourGrossW & \RelCapFourLagPfifty{} / \RelCapFourLagPninetyFive \\
8 & \RelCapEightBusy & \RelCapEightNetW{} / \RelCapEightGrossW & \RelCapEightLagPfifty{} / \RelCapEightLagPninetyFive \\
\textbf{16} & \textbf{\RelCapSixteenBusy} & \textbf{\RelCapSixteenNetW{} / \RelCapSixteenGrossW} & \textbf{\RelCapSixteenLagPfifty{} / \RelCapSixteenLagPninetyFive} \\
\bottomrule
\end{tabular}
\end{table}
